\chapter*{Preface}
There is a moment, in many proofs, when the argument suddenly becomes easy: the objects of the problem have
been arranged as vectors, and a question about them has become a question about a matrix. The number of
non-crossing paths in a grid is a determinant. The importance of a page of the web is an eigenvector. The
question whether a graph with given parameters can exist becomes the question whether certain eigenvalue
multiplicities are integers, and is answered in a few lines. A thirty-year-old conjecture about Boolean
functions falls to a single cleverly signed matrix and the interlacing of its eigenvalues. Linear algebra is
the art of finding that moment, and this book is about that art.

The first four chapters rebuild the subject on firm ground, from the point of view that pays off later:
quotients and duality, the determinant as a volume and the exterior powers behind it, eigenvalues and the
structure of a single operator, inner products and orthogonal projection. The next four chapters put the
spectrum to work. The spectral theorem comes with its relatives, the min-max principle, interlacing and the
inequalities of Weyl; the singular value decomposition gives the best approximation of a matrix by a simpler
one; the theorem of Perron and Frobenius explains why positive matrices have a dominant eigenvector, and why
random walks forget where they started; and graphs are studied through their eigenvalues, from the friendship
theorem to Moore graphs and the sensitivity conjecture. The last four chapters widen the view: symmetry, which
diagonalises whole families of matrices at once and leads to Delsarte's bounds on codes; tensors, where rank
becomes subtle and matrix multiplication becomes a research problem; the inequalities that govern positive
matrices; and linear dynamics, from recurrences and differential equations to the acceleration of learning
algorithms.

The book starts from a first course in linear algebra, the one recalled in Tome~0: vector spaces, matrices,
determinants and diagonalisation. Every theorem is proved. A few deeper results, quoted for their beauty, are
labelled as remarks with a precise reference. Each chapter opens with a question to think about, and ``Pause and
think'' boxes ask you to guess before a key result. Exercises are graded \textbf{$\star$} warm-up,
\textbf{$\star\star$} for a good afternoon and \textbf{$\star\star\star$} gems; each has a hint and a complete
solution that starts with \emph{How to find it}, because finding the idea matters more than checking the details.

Each chapter has a companion Python notebook that builds the chapter's algorithms from scratch and tests the
theorems numerically: eigenvalue algorithms, random tilings, PageRank, searches for strongly regular graphs,
image compression, Strassen's algorithm and the race between gradient descent and its accelerated cousins.
The pictures on the screen are often the fastest way to believe a theorem, and to see where its hypotheses
start to matter.
